\begin{table}[H]
\centering
\caption{Selected subfield extremes in the static metric profiles}
\label{tab:extreme_subfields}
\scriptsize
\begingroup
\setlength{\tabcolsep}{2pt}
\begin{tabular}{@{}>{\raggedright\arraybackslash}p{0.17\textwidth} >{\raggedright\arraybackslash}p{0.25\textwidth} >{\raggedright\arraybackslash}p{0.25\textwidth} >{\raggedright\arraybackslash}p{0.21\textwidth}@{}}
\hline
\textbf{Metric} & \textbf{High end} & \textbf{Low end} & \textbf{Interpretation} \\
\hline
Median Centroid Distance & Nuclear and High Energy Physics; Artificial Intelligence; Electrical and Electronic Engineering & Applied Microbiology and Biotechnology; General Energy; Research and Theory & Broadest versus most compact semantic territories. \\
Median kNN Distance & Molecular Biology; Materials Chemistry; Biomedical Engineering & Applied Microbiology and Biotechnology; Internal Medicine; General Energy & Sparsest versus tightest local neighborhoods. \\
kNN In-Degree Gini & General Materials Science; Electrochemistry; Internal Medicine & Artificial Intelligence; Computer Vision and Pattern Recognition; Cognitive Neuroscience & Most versus least hub-concentrated neighborhood structure. \\
PCA D80 & Classics; Religious studies; Electrochemistry & Business and International Management; Computational Mathematics; Human Factors and Ergonomics & Most versus least distributed PCA variance. \\
Recent Novelty Score & Computer Vision and Pattern Recognition; Signal Processing; Human Factors and Ergonomics & Pharmaceutical Science; Molecular Medicine; Animal Science and Zoology & Recent work farthest from versus closest to the early semantic core. \\
\hline
\end{tabular}
\endgroup
\par\vspace{0.15cm}
\footnotesize\textit{Note.} Extremes are computed on raw metric values among the 241 eligible subfields; table entries list the top three subfields at each end of the selected metric.
\end{table}
